% GENERATED FILE. Edit canonical lessons, then run npm run generate:books.
% canonical-source-hash: fnv1a64:7ace3eefdc6d79f4
% canonical-lessons: IT-C302-incendio, IT-C302-divampare, IT-C302-estintore, IT-C302-ustione, IT-C302-ustionarsi

\chapter{Al fuoco}
\label{ch:it-al-fuoco}
\hlchaptermodality{302}

\begin{chapteropening}
\textbf{By the end of this chapter:} \emph{I can report a fire: how it broke out, the extinguisher and the burns.}

Complete the last lesson of chapter 302: I can report a fire: how it broke out, the extinguisher and the burns.
\end{chapteropening}

\section[l'incendio]{l'incendio — the fire, the blaze}

\label{lesson:IT-C302-incendio}



\begin{quote}
Before the new one: say the Italian for the stretcher, then the Italian for the siren.
\end{quote}

\subsection*{You'll want to know: l'incendio}
\textbf{l'incendio} — \textquotedblleft{}the fire, the blaze\textquotedblright{}.

\begin{quote}
L'incendio è cominciato in cucina.
\end{quote}

\begin{quote}
\emph{The fire started in the kitchen.}
\end{quote}

A destructive fire; a fire in the fireplace is \emph{il fuoco}.

\begin{cousinweb}
\emph{incendio} is a learned borrowing from Latin \emph{incendium} \textquotedblleft{}a fire\textquotedblright{} (EtymDB), from \emph{incendere} \textquotedblleft{}to set alight\textquotedblright{}.
\end{cousinweb}

\begin{grammarlens}[title={What you've built}]
One more word for telling, explaining and describing at B1.
\end{grammarlens}

\subsection*{Your turn}
\begin{itemize}
\raggedright
  \item \emph{Say it:} \emph{l'incendio}
  \item \emph{Say it:} \emph{l'incendio}, once more
  \item \emph{Say it:} a sentence of your own with \emph{l'incendio}
  \item \emph{Recall:} say the Italian for the stretcher, then the Italian for the siren, then say \emph{l'incendio} again
\end{itemize}

\begin{tcolorbox}[breakable,colback=teal!4,colframe=teal!35!black,title={Before you move on}]
What does l'incendio mean? (\textquotedblleft{}The fire, the blaze\textquotedblright{}.) Say it once more.
\end{tcolorbox}
\section[divampare]{divampare — to flare up, to break out (of a fire)}

\label{lesson:IT-C302-divampare}



\begin{quote}
Before the new one: say the Italian for the siren, then the Italian for the fire.
\end{quote}

\subsection*{You'll want to know: divampare}
\textbf{divampare} — \textquotedblleft{}to flare up, to break out (of a fire)\textquotedblright{}.

\begin{quote}
Il fuoco è divampato in pochi minuti.
\end{quote}

\begin{quote}
\emph{The fire flared up in a few minutes.}
\end{quote}

With \emph{essere}: \emph{è divampato}.

\begin{cousinweb}
\emph{divampare} is built on \emph{vampa} \textquotedblleft{}a blaze, a flame\textquotedblright{} with \emph{di-} (EtymDB); \emph{vampa} is usually traced to Latin \emph{vapor} \textquotedblleft{}steam, heat\textquotedblright{}.
\end{cousinweb}

\begin{grammarlens}[title={What you've built}]
One more word for telling, explaining and describing at B1.
\end{grammarlens}

\subsection*{Your turn}
\begin{itemize}
\raggedright
  \item \emph{Say it:} \emph{divampare}
  \item \emph{Say it:} \emph{divampare}, once more
  \item \emph{Say it:} a sentence of your own with \emph{divampare}
  \item \emph{Recall:} say the Italian for the siren, then the Italian for the fire, then say \emph{divampare} again
\end{itemize}

\begin{tcolorbox}[breakable,colback=teal!4,colframe=teal!35!black,title={Before you move on}]
What does divampare mean? (\textquotedblleft{}To flare up, to break out (of a fire)\textquotedblright{}.) Say it once more.
\end{tcolorbox}
\section[l'estintore]{l'estintore — the fire extinguisher}

\label{lesson:IT-C302-estintore}



\begin{quote}
Before the new one: say the Italian for the fire, then the Italian for to flare up.
\end{quote}

\subsection*{You'll want to know: l'estintore}
\textbf{l'estintore} — \textquotedblleft{}the fire extinguisher\textquotedblright{}.

\begin{quote}
C'è un estintore su ogni piano.
\end{quote}

\begin{quote}
\emph{There is an extinguisher on every floor.}
\end{quote}

Built on Latin \emph{exstinguere}, ‘to put out’.

\begin{cousinweb}
\emph{estintore} is built on Latin \emph{exstinctor} \textquotedblleft{}one who puts out\textquotedblright{}, from \emph{exstinguere} \textquotedblleft{}to quench\textquotedblright{} (Lewis \& Short).
\end{cousinweb}

\begin{grammarlens}[title={What you've built}]
One more word for telling, explaining and describing at B1.
\end{grammarlens}

\subsection*{Your turn}
\begin{itemize}
\raggedright
  \item \emph{Say it:} \emph{l'estintore}
  \item \emph{Say it:} \emph{l'estintore}, once more
  \item \emph{Say it:} a sentence of your own with \emph{l'estintore}
  \item \emph{Recall:} say the Italian for the fire, then the Italian for to flare up, then say \emph{l'estintore} again
\end{itemize}

\begin{tcolorbox}[breakable,colback=teal!4,colframe=teal!35!black,title={Before you move on}]
What does l'estintore mean? (\textquotedblleft{}The fire extinguisher\textquotedblright{}.) Say it once more.
\end{tcolorbox}
\section[l'ustione]{l'ustione — the burn (on the skin)}

\label{lesson:IT-C302-ustione}



\begin{quote}
Before the new one: say the Italian for to flare up, then the Italian for the fire extinguisher.
\end{quote}

\subsection*{You'll want to know: l'ustione}
\textbf{l'ustione} — \textquotedblleft{}the burn (on the skin)\textquotedblright{}.

\begin{quote}
Ha un'ustione sul braccio.
\end{quote}

\begin{quote}
\emph{He has a burn on his arm.}
\end{quote}

A doctor's word for a burn on the skin.

\begin{cousinweb}
\emph{ustione} is a learned borrowing from Latin \emph{ustiō} \textquotedblleft{}a burning\textquotedblright{}, from \emph{ūrere} \textquotedblleft{}to burn\textquotedblright{} (Lewis \& Short).
\end{cousinweb}

\begin{grammarlens}[title={What you've built}]
One more word for telling, explaining and describing at B1.
\end{grammarlens}

\subsection*{Your turn}
\begin{itemize}
\raggedright
  \item \emph{Say it:} \emph{l'ustione}
  \item \emph{Say it:} \emph{l'ustione}, once more
  \item \emph{Say it:} a sentence of your own with \emph{l'ustione}
  \item \emph{Recall:} say the Italian for to flare up, then the Italian for the fire extinguisher, then say \emph{l'ustione} again
\end{itemize}

\begin{tcolorbox}[breakable,colback=teal!4,colframe=teal!35!black,title={Before you move on}]
What does l'ustione mean? (\textquotedblleft{}The burn (on the skin)\textquotedblright{}.) Say it once more.
\end{tcolorbox}
\section[ustionarsi]{ustionarsi — to get burned, to burn oneself}

\label{lesson:IT-C302-ustionarsi}



\begin{quote}
Before the new one: say the Italian for the fire extinguisher, then the Italian for the burn.
\end{quote}

\subsection*{You'll want to know: ustionarsi}
\textbf{ustionarsi} — \textquotedblleft{}to get burned, to burn oneself\textquotedblright{}.

\begin{quote}
Mi sono ustionato la mano.
\end{quote}

\begin{quote}
\emph{I burned my hand.}
\end{quote}

For burns on the skin: \emph{si è ustionato le dita}.

\begin{cousinweb}
\emph{ustionarsi} is built on \emph{ustione} \textquotedblleft{}burn\textquotedblright{}, from Latin \emph{ustiō} \textquotedblleft{}a burning\textquotedblright{}, from \emph{ūrere} \textquotedblleft{}to burn\textquotedblright{} (Lewis \& Short).
\end{cousinweb}

\begin{grammarlens}[title={What you've built}]
One more word for telling, explaining and describing at B1.
\end{grammarlens}

\subsection*{Your turn}
\begin{itemize}
\raggedright
  \item \emph{Say it:} \emph{ustionarsi}
  \item \emph{Say it:} \emph{ustionarsi}, once more
  \item \emph{Say it:} a sentence of your own with \emph{ustionarsi}
  \item \emph{Recall:} say the Italian for the fire extinguisher, then the Italian for the burn, then say \emph{ustionarsi} again
\end{itemize}

\begin{tcolorbox}[breakable,colback=teal!4,colframe=teal!35!black,title={Before you move on}]
What does ustionarsi mean? (\textquotedblleft{}To get burned, to burn oneself\textquotedblright{}.) Say it once more.
\end{tcolorbox}
